\section{Lessons and Path Forward}
\label{sec:lessons}

Three observations from this work generalize beyond the MRA-NN setting
and may inform future attempts to couple ML models with iterative
scientific solvers.

Input signal quality determines the ceiling of any downstream model,
independent of architecture capacity. Our parent model doubled the
input dimension (2816 vs.\ 1792 features) and added 34\% more
parameters (4.1M vs.\ 3.0M), yet reproduced the per-level MSE ratios
of the base model almost exactly. The single-task ablation confirmed
that multi-task interference was a real training pathology, but
removing it only brought the model to parity with $\rho_0$, not beyond
it. The fundamental limitation is that the promolecular density and
nuclear potential carry no electron-electron correlation signal. No
amount of nonlinear capacity can extract exchange-correlation physics
from inputs that do not encode it.

Fine-scale accuracy does not imply solver acceleration. The level-clamped
model achieved MSE ratios of 0.41--0.98$\times$ baseline at levels
10--14, a measurable and statistically significant improvement over
$\rho_0$. Yet SCF injection tests produced identical iteration counts
across all 14 molecule-threshold combinations (7 molecules, 2
thresholds). The explanation is structural: MADNESS
protocols~\cite{harrison2016madness} refine from coarse to fine,
and iteration count is determined by the coarsest protocol stage
where the initial density deviates from the converged solution.
Improvements at spatial scales finer than the first protocol
threshold provide no leverage on convergence. Any ML model targeting
solver acceleration must improve the density at the resolution the
solver actually queries first, not at the resolution where data is
most abundant.

Multi-task interference in uncertainty-weighted losses can silently
undermine the primary objective. Kendall uncertainty
weighting~\cite{kendall2018multitask} allowed the learned scale
parameter $\sigma_{\rho_s}$ to grow without bound, progressively
zeroing the density gradient while the classification head consumed
the entire learning signal. The resulting model was 2.99$\times$
worse than $\rho_0$ on density prediction, yet the aggregate loss
decreased monotonically, masking the failure. A subsequent
configuration flag (\texttt{refine\_focused}) appeared to fix this
by switching the checkpoint selection metric to refinement F1, but it
left the loss function unchanged; the density head still trained under
the corrupted gradient. Refine-only training with pure focal loss
achieved F1\,=\,0.860, a gain of only 0.003 over the multi-task
formulation (F1\,=\,0.857). The broader warning is that multi-task
losses with learned weighting parameters require per-head monitoring
of both loss and scale; aggregate convergence is not a reliable
diagnostic.

\subsection{Path Forward: Gaussian-Basis Density as ML Input}
\label{sec:path_forward}

The negative results above point to a concrete next step. If the
bottleneck is input signal quality, the remedy is to supply an input
that already contains exchange-correlation physics. A natural candidate
is the converged density from a Gaussian-basis DFT calculation (e.g.,
Dalton), projected onto the MRA grid. Preliminary experiments by
collaborators show that using the Dalton density directly as the
MADNESS starting guess reduces SCF iteration counts, confirming that
the missing signal is the core issue.

This reframes the ML problem from predicting a large correction
(promolecular $\rho_0$ $\to$ converged MRA $\rho$) to predicting a
small residual (Gaussian-basis $\rho_{\text{GTO}}$ $\to$ converged MRA
$\rho$). The residual is dominated by basis-set truncation artifacts
near nuclei and in diffuse tails, both of which are spatially localized
and well-suited to the local halo-feature architecture already
developed. We plan to construct training data by pairing Dalton
densities with converged MADNESS densities for the same molecules and
to retrain the single-task model on the resulting residuals. No
results from this approach are available yet, and the magnitude of the
achievable speedup remains an open question.
